\documentclass[11pt,a4paper]{article}
\usepackage[T1]{fontenc}
\usepackage[utf8]{inputenc}
\usepackage[main=italian,provide=*]{babel}
\usepackage{amsmath,amssymb}
\usepackage{geometry}
\geometry{margin=2.5cm}
\usepackage{booktabs}
\usepackage{seqsplit}
\usepackage{hyperref}
\hypersetup{colorlinks=true,linkcolor=blue,citecolor=blue,urlcolor=blue}

\newcommand{\ket}[1]{\lvert #1 \rangle}

\title{Correlazioni dinamiche — trimero a catena aperta}
\author{Note di lavoro — Tesi triennale, Samuele Schiavi\\ Relatore: Prof.\ Alessandro Chiesa}
\date{}

\begin{document}
\maketitle

\section*{Introduzione}

Questo documento raccoglie i file della fase di \emph{correlazioni dinamiche} (circuito con
l'ancilla, Hadamard test) per il trimero a catena aperta — mirror diretto, per la topologia a due
legami, di quanto già completato per l'anello
(\texttt{\seqsplit{trimero\_anello\_correlazioni.tex}}). A differenza dell'anello, questa fase
lavora fin dall'inizio su \textbf{due} punti di lavoro, non uno solo (VQE-DM,
$J{=}1,b{=}b_c{=}3.0,D{=}0.15$; S0 Trotter, $J{=}1,b{=}3.0073414,D{=}0.3$), tenuti entrambi per
confronto perché la scelta definitiva resta aperta col relatore. L'ordine seguito è stato lo
stesso dell'anello: prima l'analisi di simmetria (verificata numericamente e sottoposta a tre
giri di audit indipendenti, con correzioni reali a ogni giro), poi il circuito.

Per ciascun file: cosa fa, le funzioni interne, come si usa, le verifiche effettuate, e cosa se ne
può concludere.

\section*{\texttt{\seqsplit{simmetrie\_correlatori\_trimero\_catena.tex}}}

Documento di teoria, primo passo della fase: analisi di simmetria dell'Hamiltoniana completa
(scambio + campo + DM) applicata ai correlatori
$C_{ij}^{\alpha\beta}(t)=\langle\psi_0|\sigma_i^\alpha(t)\sigma_j^\beta(0)|\psi_0\rangle$, prima di
qualunque implementazione di circuito.

\textbf{Contenuto.} La simmetria residua è $P_{13}$ (riflessione siti $1,3$, sito $2$ fisso) --
uno \textbf{scambio puro}, senza rotazioni aggiuntive, a differenza dell'anello che richiede
$U_\text{anello}=\mathrm{SWAP}_{12}\cdot(R_z(\pi))^{\otimes3}$. Conseguenza qualitativa: i
coefficienti $\lambda_\alpha\equiv+1$ per ogni $\alpha$ (contro $\eta_x=\eta_y=-1,\eta_z=+1$
dell'anello), quindi $P_{13}$ produce \textbf{solo uguaglianze} fra correlatori distinti
($C_{11}\equiv C_{33}$, $C_{12}\equiv C_{32}$, $C_{21}\equiv C_{23}$, $C_{13}\equiv C_{31}$),
\textbf{mai zeri}: il sito fisso (2) non è vincolato, a differenza del sito fisso (3) dell'anello
(4 zeri strutturali). Time-reversal $K$: $24$ zeri a $t=0$ per $i\neq j$ (stessa dimostrazione
dell'anello) più $6$ zeri aggiuntivi per $i=j$ ($C_{ii}^{xz},C_{ii}^{zx}$), da un argomento
distinto ($\langle\sigma_i^y\rangle=0$ per stato reale). A $D=0$: regola di selezione su $M$, $36$
zeri per ogni $t$, distrutti tutti dal DM (contro $56\to4$ dell'anello).

\textbf{Verifiche.} Tre giri di audit indipendenti, ciascuno con correzioni reali: (1) errore di
conteggio nel Cor.\ delle relazioni $P_{13}$ (72 correlatori in 36 coppie, non 36 correlatori),
più due lacune (caso $D=0$ mai trattato, degenerazione del punto di lavoro a $D=0$ mai
dichiarata); (2) collisione di notazione ($\eta_\alpha$ usato con due significati diversi in due
lemmi, rinominato $\lambda_\alpha$ per $P_{13}$); (3) un claim comparativo verso l'anello errato
($56\to4$ per l'anello, non "restava 4"). Scan completo delle 81 combinazioni, due metodi
indipendenti (esponenziale di matrice diretto e formula spettrale), confrontati direttamente fra
loro (accordo $5.5\times10^{-15}$).

\textbf{Conclusione.} Nessuno zero strutturale per ogni $t\neq0$ (0/81, contro i 4 dell'anello);
$30/81$ nulle a $t=0$; $45$ valori indipendenti determinano tutte le $81$ combinazioni. Punto di
lavoro non degenere solo grazie al DM (verificato: degenere a $D=0$).

\section*{\texttt{\seqsplit{circuito\_correlazioni\_trimero\_catena.py}}}

Modulo di produzione principale: costruisce il circuito Hadamard test a 4 qubit (3 di registro +
1 ancilla) per misurare $C_{ij}^{\alpha\beta}(t)$, mirror diretto del circuito già validato per
l'anello, riusando \texttt{\seqsplit{trotter\_trimero\_catena.py}} come blocco $U(t)$.

\textbf{Funzioni interne.}
\begin{itemize}
\item \texttt{ANCILLA = 3} — indice del qubit ancilla.
\item \texttt{ground\_state(J,\ b,\ D)} — stato fondamentale esatto (fase reale fissata).
\item \texttt{build\_correlator\_circuit(...)} — il circuito completo (stessi argomenti
dell'anello più \texttt{alternate}, propagato a \texttt{trotter\_circuit}, opzione assente nel
modulo dell'anello).
\item \texttt{ancilla\_z\_expectation(qc)}, \texttt{correlator\_from\_circuit(...)} — lettura e
combinazione Re/Im, identiche nella struttura all'anello.
\end{itemize}

\textbf{Verifiche.} Test mirato con \textbf{stato iniziale asimmetrico} sotto $P_{13}$
($\|P_{13}\psi-\psi\|=1.21$): con il ground state (autostato di $P_{13}$) un mapping
sito$\leftrightarrow$qubit invertito sarebbe stato invisibile a ogni altro test. Errore su tutte
le 81 combinazioni scala $O(1/N)$ (rapporto $2.00$ raddoppiando $N$, da $N{=}50$ a $N{=}800$):
puro Trotter, nessun bias residuo — mapping verificato corretto in condizioni che ne avrebbero
rivelato l'errore.

\textbf{Conclusione.} Modulo maturo, usato senza modifiche da tutti gli script successivi.

\section*{\texttt{\seqsplit{vqe\_w2qC\_k2\_trimero\_catena.py}}}

Modulo minimo per l'ansatz \texttt{W-2qC.K2} (\texttt{pma\_2q\_trimer\_cyclic(2, w\_block)}, 10
parametri, 2 cicli), mirror di \texttt{vqe\_w2q6\_trimero\_anello.py}. A differenza dell'anello,
ottimizza \textbf{entrambi} i punti di lavoro nella stessa esecuzione (\texttt{main()}).

\textbf{Risultato.} $60$ restart (COBYLA + polish L-BFGS-B) per punto: VQE-DM
$\mathcal F=1.000000000000$; S0 $\mathcal F=0.999999999999992$. Nella sessione di sviluppo,
$5$--$15$ restart erano già sufficienti; $60$ adottati per lo stesso motivo di cautela già
documentato in \texttt{\seqsplit{trimero\_catena\_vqe\_dm.tex}} (vicinanza a un level crossing a
$D=0$). Parametri ricaricati e ricalcolati indipendentemente (energia, fidelity, norma dello
stato): coincidenti coi valori salvati.

\textbf{Conclusione.} Due file \texttt{.npz}, uno per punto, non intercambiabili.

\section*{\texttt{\seqsplit{validate\_circuito\_correlazioni\_catena.py}}}

Validazione indipendente: riferimento classico esatto per esponenziale di matrice diretto (non
formula spettrale), mirror di \texttt{validate\_circuito\_correlazioni.py}. Esegue
\textbf{automaticamente entrambi} i punti di lavoro.

\textbf{Verifiche.} Convergenza Trotter confermata ($O(1/N)$, rapporto $\to2.0$); zeri predetti a
$t=0$ confermati a precisione macchina, sia sull'esatto sia sul circuito; relazione $P_{13}$
($C_{11}^{xx}\equiv C_{33}^{xx}$) rispettata entro l'errore di Trotter, in entrambi i punti.

\section*{\texttt{\seqsplit{validate\_vqe\_circuito\_correlazioni\_catena.py}}}

Chiude la pipeline VQE$\to$correlazioni, mirror di
\texttt{validate\_vqe\_circuito\_correlazioni.py}. Confronto sistematico delle 81 combinazioni fra
preparazione VQE reale e preparazione esatta, su entrambi i punti.

\textbf{Risultato.} Errore massimo $5.9\times10^{-8}$ (VQE-DM) e $9.8\times10^{-8}$ (S0), entrambi
coerenti con $\sqrt{1-\mathcal F}$ — stesso comportamento dell'anello.

\section*{\texttt{\seqsplit{scan81\_trimero\_catena.py}} + \texttt{\seqsplit{generate\_figure\_scan81\_catena.py}}}

Estensione a tutte e 81 le combinazioni via circuito effettivo (statevector e shot finiti, 8192
shots), mirror di \texttt{scan81\_trimero\_anello.py}. A differenza dell'anello, prende un
argomento da riga di comando (\texttt{vqedm} o \texttt{s0}) e va lanciato una volta per punto.

\textbf{Risultati.}
\begin{center}
\begin{tabular}{lcc}
\toprule
& VQE-DM ($D{=}0.15$) & S0 ($D{=}0.3$) \\
\midrule
$|C|$ minimo statevector & $0.1008$ & $0.0168$ \\
errore Trotter, media/max & $1.76\times10^{-3}$/$3.34\times10^{-3}$ & $1.72\times10^{-3}$/$3.97\times10^{-3}$\\
errore statistico, media/max & $1.39\times10^{-2}$/$3.85\times10^{-2}$ & $1.36\times10^{-2}$/$3.90\times10^{-2}$\\
\bottomrule
\end{tabular}
\end{center}
Nessuna delle 81 combinazioni collassa a zero, in nessuno dei due punti — coerente col documento
di simmetria (0 zeri predetti, contro i 4 dell'anello).

\section*{\texttt{\seqsplit{validate\_shot\_noise\_trimero\_catena.py}} + \texttt{\seqsplit{generate\_figure\_shotnoise\_catena.py}}}

Separazione errore Trotter/statistico, mirror di
\texttt{validate\_shot\_noise\_trimero\_anello.py}. Un punto per lancio (argomento da riga di
comando).

\textbf{Risultato.} Rapporto statistico/Trotter $\sim10\times$ a VQE-DM, $\sim115\times$ a S0
(su $C_{11}^{zz}(t{=}1.3)$): lo statistico domina in entrambi, con margine molto diverso —
specifico del correlatore e del punto scelti, non un'anomalia. Convergenza $1/\sqrt{N_\text{shots}}$
confermata su 7 valori di shot $\times$ 40 ripetizioni, in entrambi i punti.

\section*{\texttt{\seqsplit{circuito\_correlazioni\_trimero\_catena\_spiegato.tex}}}

Documento pedagogico principale: spiegazione passo per passo del circuito, mirror di
\texttt{\seqsplit{circuito\_correlazioni\_trimero\_anello\_spiegato.tex}} -- non una copia: la
Sez.~2 è ribaltata rispetto all'anello (mostra \emph{perché} il sito fisso \textbf{non} produce
zeri, con figura di contrasto diretto $C_{22}^{xz}(t)$), e la Sez.~6 aggiunge una novità assente
nell'anello (l'opzione \texttt{alternate}, guadagno non monotono in $D$ e $N$).

\textbf{Contenuto specifico della catena.} Verifica dedicata sul mapping sito$\leftrightarrow$qubit
con stato asimmetrico (Sez.~8, box "Verifica mirata"); tabella "$U$ controllato" con valori
effettivamente calcolati (non placeholder, corretto durante la stesura); scan a 81 combinazioni su
entrambi i punti di lavoro.

\textbf{Figure prodotte} (da \texttt{\seqsplit{generate\_circuit\_fig\_trimero\_catena.py}},
\texttt{\seqsplit{generate\_figures\_circuito\_trimero\_catena.py}}): diagramma del circuito
(Re/Im), convergenza Trotter, validazione continua $C_{11}^{yy}(t)$, contrasto "nessuno zero sul
sito 2" ($C_{22}^{xz}(t)$).

\textbf{Verifiche.} Compilato due volte, zero errori, un solo underfull cosmetico residuo (stessa
soglia di tolleranza già accettata nei documenti dell'anello). 10 pagine, tutte verificate
visivamente pagina per pagina.

\section*{Quattro notebook della fase}

Mirror per la catena della struttura a quattro notebook già usata per l'anello, eseguiti per
intero (non solo listati di codice), tutta la logica pesante riusata per import dai moduli sopra.

\textbf{\texttt{\seqsplit{correlazioni\_trimero\_catena\_simmetria.ipynb}}} (18 celle) — mirror
computazionale del documento di teoria: $P_{13}$ commuta senza rotazioni aggiuntive, $\lambda_\alpha
\equiv+1$, nessuno zero forzato sul sito fisso, i due corollari di time-reversal (30 zeri a
$t=0$), la regola $D=0$ (36 zeri). Ogni numero verificato coincidere con quello del \texttt{.tex}.

\textbf{\texttt{\seqsplit{correlazioni\_trimero\_catena\_esplorazione.ipynb}}} (16 celle) —
circuito con l'ancilla disegnato esplicitamente, validazione su 4 casi, convergenza Trotter, shot
noise su un caso singolo.

\textbf{\texttt{\seqsplit{circuito\_correlazioni\_trimero\_catena\_tutte.ipynb}}} (25 celle) —
scan completo delle 81 combinazioni con preparazione esatta: heatmap, griglia $9\times9$ nel
tempo (Re \emph{e} Im — corretto durante la sessione dopo una domanda mirata sull'assenza della
parte immaginaria), selettore, analisi spettrale, verifica esplicita che nessun correlatore sia
quasi-costante (escursione minima $0.196$, contro $\sim0.0005$ di $C_{33}^{zz}$ nell'anello) --
anch'essa aggiunta dopo la stessa domanda, non presente nella prima stesura.

\textbf{\texttt{\seqsplit{circuito\_correlazioni\_trimero\_catena\_vqe.ipynb}}} (22 celle) —
stessa analisi con preparazione VQE reale. Residuo medio $1.89\times10^{-8}$, identico a quello
dello script standalone corrispondente: buon controllo incrociato fra le due strade.

\textbf{Verifiche.} Tutti e quattro eseguiti senza errori; tutti i numeri stampati verificati
coincidere con quelli degli script standalone dove sovrapposti.

\textbf{Conclusione.} Mirror computazionale completo; una correzione reale emersa durante
l'ispezione (parte immaginaria mancante nella griglia temporale, assenza di verifica esplicita
sui correlatori quasi-costanti), non solo confermata ma con causa e correzione dichiarate.

\section*{\texttt{\seqsplit{manuale\_uso\_correlazioni\_trimero\_catena\_\{statevector,vqe\}.tex}}}

Manuali d'uso, mirror dei corrispondenti dell'anello. Differenza strutturale dichiarata in
entrambi: la catena lavora su \textbf{due} punti di lavoro, non uno; gli script che li supportano
prendono un argomento da riga di comando (\texttt{vqedm}/\texttt{s0}) oppure eseguono entrambi
automaticamente, a seconda dello script -- distinzione esplicita nel manuale, verificata contro
il comportamento reale di ciascuno script prima di scriverla.

\textbf{Verifiche.} Compilati entrambi, zero errori, alcuni overfull cosmetici residui (righe di
codice/nomi file lunghi in \texttt{verbatim}/\texttt{texttt}), verificati non compromettere la
leggibilità pagina per pagina.

\section*{Conclusione generale}

La fase di correlazioni dinamiche per la catena aperta è arrivata a copertura completa: teoria di
simmetria derivata e sottoposta a tre giri di audit indipendenti (con correzioni reali a ogni
giro, non solo conferme), circuito costruito e validato su entrambi i punti di lavoro in uso nel
progetto, pipeline VQE$\to$correlazioni chiusa, scan sistematico delle 81 combinazioni (nessuno
zero, coerente con la teoria), quattro notebook eseguiti, manuali d'uso completi. Il risultato
qualitativo centrale -- la catena non eredita gli zeri strutturali dell'anello, perché la sua
simmetria residua ($P_{13}$) è un puro scambio privo di segno -- è stato verificato con lo stesso
livello di rigore già applicato all'anello, incluso un controllo dedicato a escludere che la
simmetria stessa potesse mascherare un errore di implementazione. Aperti per il seguito: scelta
definitiva del punto di lavoro con il relatore; correlatori dinamici già chiusi per entrambe le
topologie di $N=3$ -- unico filone rimasto è la Parte 2 della tesi (rumore, Kraus/Lindblad).

\end{document}
